\documentclass[11pt]{article}
\usepackage[margin=27mm]{geometry}
\usepackage{amsmath,amssymb,amsthm,booktabs,hyperref}
\newtheorem{theorem}{Theorem}
\newtheorem{lemma}[theorem]{Lemma}
\newtheorem{proposition}[theorem]{Proposition}
\newcommand{\F}{\mathbb F}
\newcommand{\PSL}{\operatorname{PSL}}
\newcommand{\PGL}{\operatorname{PGL}}
\newcommand{\PG}{\operatorname{PG}}
\newcommand{\Aut}{\operatorname{Aut}}
\title{Subfield planes and unordered pairs of parabolics}
\author{Human Compression Project}
\date{}
\begin{document}\maketitle
\begin{abstract}
We count two families from the fixed Part II classification of rank two
geometries for $\PSL(2,q)$. In characteristic two, planes through a fixed
subfield line give the parameter space; Frobenius gives the equivalence
relation. This recovers the numerical formulas for the $A_5,E_{16}:3$
family and extends its parameter description to the previously omitted
subfield covers. In characteristic three, unordered pairs in an additive
quotient give the family with two $A_4$ parabolics. This halves its printed
conjugacy count under the source's stated convention, while preserving its
automorphism count. Classical subgroup and automorphism theorems remain
explicit premises.
\end{abstract}

\section{Local actions and the counting convention}
Let $G$ be a finite simple projective group. For distinct proper subgroups
$H,Q$, put $B=H\cap Q$. The rank two coset geometry has the cosets of
$H,Q$ as its two types, with nonempty intersection defining incidence.
The two local actions are those of $H,Q$ on their cosets of $B$.
When both are two-transitive and one parabolic is maximal in $G$, these
subgroups generate $G$ and give a connected, flag-transitive, residually
weakly primitive geometry. We use the proof of this assertion in the
previously accepted projective-line manuscript; it depends only on coset
actions and maximality.

The fixed Part II source, Section 2 on physical page 4, defines conjugacy
by equality of the \emph{sets} $\{H^g,Q^g\}$ and $\{H',Q'\}$.
Thus interchanging the two types is allowed. We call this the unordered
convention. We will give ordered counts separately when they differ.
Equivalence under $\Aut(G)$ uses the same convention.

We retain the following classical facts used by the source: the natural
subfield embeddings and their maximality at prime extension degree;
the conjugacy of those embeddings in the cases specified below;
$\Aut(\PSL(2,p^f))=\operatorname{P\Gamma L}(2,p^f)$ for the ambient
groups used here; and the usual projective orders. These are not assumed
consequences of the geometry classification under audit. Every normalizer
needed for the new counts is proved below.

\section{Planes through a subfield line}
Let $K=\F_s\subset F=\F_{s^m}$, with $s\ge4$ even and $m$ prime.
Write $U_T=\{x\mapsto x+t:t\in T\}$ and
$C=\{x\mapsto ax:a\in K^*\}$. Set
\[
 G=\PSL(2,F)=\PGL(2,F),\quad H=\PSL(2,K),\quad B=U_K C.
\]
For a two-dimensional $K$-subspace $T$ containing $K$, put $Q_T=U_TC$.
The earlier accepted subfield-plane proof establishes
\[
 H\cap Q_T=B,\qquad |H|=s(s^2-1),\quad |Q_T|=s^2(s-1),
 \quad |B|=s(s-1),
\]
and proves all geometry hypotheses. Its argument is a scalar-subspace
correspondence: groups between $U_KC$ and $U_TC$ are the groups $U_WC$
for $K$-subspaces $K\subseteq W\subseteq T$. A cover has quotient
dimension one and induces $\operatorname{AGL}(1,K)$ on its local residue.
No subfield of degree two is needed.

\begin{lemma}\label{normalizer}
The normalizer of $B$ in $G$ is $B$. For each conjugate of $B$, there
is exactly one conjugate of $H$ containing it. The conjugacy classes of
the displayed geometries are consequently indexed by
$\PG(m-2,s)$, the lines of $F/K$.
\end{lemma}
\begin{proof}
The translation subgroup $U_K$ is the unique Sylow two-subgroup of $B$,
and its common projective fixed point is infinity. A normalizer of $B$
is therefore affine, say $x\mapsto ax+b$. Normalizing $U_K$ requires
$aK=K$, hence $a\in K^*$. Conjugating any nonidentity scalar in $C$
then requires $(1-c)b\in K$, so $b\in K$. This proves the normalizer
assertion, including its converse.

The subfield groups $H$ form one $G$-class in this even-characteristic
case. Their point stabilizers form one class inside $H$. If $H^g$
contains $B$, conjugate a point stabilizer inside $H$ before applying
$g$; the resulting transporter normalizes $B$. It belongs to $B$ by
the first paragraph, so $H^g=H$.

To compare two displayed pairs, the unequal parabolic orders first force
$H$ to be preserved. Their intersections then force $B$ to be preserved.
The transporter is in $N_G(B)=B$. Its translations lie in $K\subset T$
and its scalars in $K^*$, so it normalizes each $Q_T$. Hence two planes
give conjugate pairs exactly when they are equal. The planes containing
$K$ correspond to the lines of $F/K$, as asserted.
\end{proof}

\begin{proposition}\label{general}
There are exactly $(s^{m-1}-1)/(s-1)$ conjugacy classes in this family.
Their $\Aut(G)$-classes are the orbits of the absolute Frobenius
$x\mapsto x^2$ on $\PG(m-2,s)$.
\end{proposition}
\begin{proof}
The conjugacy count is Lemma~\ref{normalizer}. Every field automorphism
preserves the standard $H$ and $B$. Write an arbitrary ambient
automorphism as an inner automorphism followed by a field automorphism.
If it transports one of our pairs to another, its inner part must
normalize $B$ and therefore belongs to $B$. This part acts trivially
on the planes. Conversely every Frobenius power gives the stated
transport. This proves the orbit description, with no assertion that
the field automorphism action is free.
\end{proof}

\begin{lemma}\label{affine}
An affine subgroup $Q$ containing $B$ with $B$ maximal in $Q$ is one of
the groups $Q_T$ above. In particular every such affine extension has
a two-transitive local action.
\end{lemma}
\begin{proof}
Write $V=Q\cap U_F$ for its translations and $D$ for its image in
$F^*$. Both $V$ and $K$ are stable under $K^*$, so $V$ is a
$K$-space. The subgroup $U_VC$ lies between $B$ and $Q$. If $V>K$,
maximality forces $Q=U_VC$, and the scalar correspondence forces
$\dim_K(V/K)=1$. If $V=K$, an element $x\mapsto ax+b$ of $Q$
normalizes $U_K$, so $a\in K^*$. Multiplying by its scalar inverse
in $C$ leaves translation by $b\in V=K$. Thus $Q=B$, contrary to
the proper extension. This excludes the second case.
\end{proof}

\section{The $A_5,E_{16}:3$ formulas}
Specialize to $K=\F_4$, $F=\F_{4^r}$ with $r$ prime. The numbers
in this section are the ones asserted by Lemmas 5.3--5.4 and the
corresponding rows of Table 1 in the fixed source.

\begin{lemma}\label{subgroups}
For odd prime $r$, the groups isomorphic to $E_{16}:3$ in $G$ form
$(4^{r-1}-1)/15$ conjugacy classes.
\end{lemma}
\begin{proof}
Such a group has characteristic normal subgroup $E_{16}$, whose
nonidentity elements are unipotent and have a unique common fixed
projective point. Put this point at infinity. Its complement of order
three has a second fixed point; a translation puts that point at zero.
The group is now $U_TC$ for a $K$-plane $T\subset F$. Conversely
each such plane supplies the stated group. Affine conjugation identifies
exactly the scalar multiples of the planes: translations can change a
chosen complement but cannot change the translation subgroup.

The multiplier ring $E_T=\{a\in F:aT\subseteq T\}$ is a field.
Closure under addition and multiplication is immediate; for nonzero
$a\in E_T$, multiplication is a bijection of the finite set $T$,
so $a^{-1}\in E_T$. It contains $K$, and $[E_T:K]$ divides both
$\dim_K T=2$ and $[F:K]=r$. Thus $E_T=K$. Every scalar orbit has
length $(4^r-1)/3$. There are
\[
 {r\brack 2}_4=\frac{(4^r-1)(4^r-4)}{(16-1)(16-4)}
\]
planes, by counting ordered independent pairs and dividing by the number
of ordered bases of a plane. Division by the orbit length gives
$(4^{r-1}-1)/15$.
\end{proof}

\begin{theorem}\label{four}
For $r=2$, the family with parabolics $A_5,E_{16}:3$ has one class
under each equivalence. For $r=3$ it has five conjugacy classes and two
$\Aut(G)$-classes. For prime $r>3$ the respective counts are
\[
 N=\frac{4^{r-1}-1}{3},\qquad
 \frac{N+(2^{r-1}-1)}{2r}
 =\frac{2(4^{r-2}-1)+3\cdot2^{r-2}}{3r}.
\]
\end{theorem}
\begin{proof}
The planes through $K$ form $\PG(r-2,4)$, of size $N$; when $r=2$
there is just the plane $F$. Put $V=F/K$, and let $\sigma$ be squaring
on this quotient. We must count its orbits, by Proposition~\ref{general}.
It has order dividing $2r$. Put $\tau=\sigma^2$, the $K$-linear
Frobenius of order $r$, and $\beta=\sigma^r$, an involution semilinear
over $K/\F_2$.

First suppose $r>3$. If a line is fixed by $\tau$, choose a nonzero
vector $v$ on it with $\tau v=av$, $a\in K^*$. Iterating gives
$a^r=1$, so $a=1$ because $r$ is prime and different from three.
Writing $v=x+K$ gives $x^4-x=c\in K$. The relative trace from $F$
to $K$ annihilates the left side and takes $c$ to $rc=c$, so $c=0$.
Then $x\in K$, a contradiction. Every nonidentity power of $\tau$
has the same fixed lines. An odd power of $\sigma$ other than $\beta$
generates the cyclic group of order $2r$, so it has no fixed line either.

It remains to count the lines fixed by $\beta$. If $\beta(v)=av$,
then $a\in K^*$ and one may replace $v$ by $a^{-1}v$ to make it
fixed: $\beta(a^{-1})a=a^{-1}$ since $a^3=1$. Each stable $K$-line
therefore contains exactly one nonzero fixed vector. The fixed quotient
space $V^\beta$ is naturally $\F_{2^r}/\F_2$. Indeed, for a fixed
class $x+K$, $c=\beta x-x\in K$ satisfies $\beta c=c$, hence
$c\in\F_2$. The map $k\mapsto k^2-k$ from $K$ onto $\F_2$
provides $k$ with $\beta k-k=c$; subtracting $k$ gives a representative
fixed by $\beta$. Its ambiguity is exactly $\F_2$. Thus $\beta$
fixes $2^{r-1}-1$ projective points. Burnside's formula gives the claimed
average over $2r$ powers.

For $r=3$, $V$ has dimension two over $K$. The trace argument still
shows that $\tau$ has no fixed nonzero vector. Since $\tau^3=1$,
it is semisimple and its possible eigenvalues are the two nonidentity
cube roots in $K$. The commuting semilinear map $\beta$ squares
eigenvalues and is invertible, so both eigenspaces occur and have the
same dimension. Each therefore has dimension one. These are precisely
the two invariant lines, and $\beta$ interchanges them. The odd
generators $\sigma,\sigma^5$ also
interchange them and fix no line. The preceding fixed-quotient argument
gives three lines fixed by $\beta$. The six fixed-point counts are
$5,0,2,3,2,0$, whose average is two.
\end{proof}

\section{Two $A_4$ parabolics in characteristic three}
Let $q=3^m$, where $m$ is an odd prime, and put
$G=\PSL(2,q)$, $H=\PSL(2,3)=A_4$, $B=U_{\F_3}\cong C_3$.
The classical subfield theorem makes $H$ maximal in $G$; its conjugates
form one class. Put $V=\F_q/\F_3$, regarded as an additive group.

\begin{lemma}\label{torsor}
The conjugates of $H$ containing $B$ form a regular $V$-set under
translations. Distinct such subgroups intersect exactly in $B$.
Conjugacy classes of ordered pairs of them are indexed by $V\setminus0$;
unordered classes are indexed by $(V\setminus0)/\{v\sim-v\}$.
\end{lemma}
\begin{proof}
Any normalizer of $B$ fixes its unique projective fixed point and is
affine. Its slope must lie in $\F_3^*=\{1,-1\}$. An affine
transformation belongs to $\PSL(2,q)$ precisely when its slope is a
square in $\F_q$. Since $m$ is odd, $q\equiv3\pmod4$ and $-1$
is not a square. Therefore $N_G(B)=U_{\F_q}$.

Every copy of $B$ inside $H$ is conjugate under $H$, and $N_H(B)=B$.
Thus the copies of $H$ containing the fixed $B$ are exactly its
translates $H_t=u_t H u_t^{-1}$, with equality precisely when
$t-t'\in\F_3$. Here we also use $N_G(H)=H$, which follows from
maximality and simplicity. Translation acts regularly on the index set
$V$. Since $B$ is maximal in $A_4$, any two distinct $H_t$ meet in
$B$.

Every eligible pair of $A_4$ subgroups with intersection of order three
can be transported to this set, since its intersection is a conjugate of
$B$. Any transporter between two such pairs must normalize their
intersection. Hence it acts as translation on $V$. Translating an
ordered pair $(a,b)$ to $(0,b-a)$ gives precisely its nonzero difference.
Interchanging the entries changes this difference to its negative.
These are all identifications, proving both parameterizations.
\end{proof}

\begin{theorem}\label{three}
Under the unordered convention of the fixed source, the number of
conjugacy classes of these geometries is
\[
 \frac{3^{m-1}-1}{2}.
\]
The number of $\Aut(G)$-classes is two if $m=3$, and
$(3^{m-1}-1)/(2m)$ otherwise. The ordered conjugacy count is
$3^{m-1}-1$.
\end{theorem}
\begin{proof}
Distinct pairs give the required geometries: their intersection is $C_3$,
each local $A_4$ action on four cosets is two-transitive, and maximality
of either $A_4$ gives connectedness and the residual primitivity conditions.
In characteristic three a nonzero vector is different from its negative.
Lemma~\ref{torsor} therefore proves the two conjugacy counts.

In $\operatorname{P\Gamma L}(2,q)$ the normalizer of $B$ consists
exactly of affine maps with slope $\pm1$, followed by field automorphisms.
The standard $H$ is preserved by sign and by every field automorphism.
It follows that the additional action on differences is generated by
$v\mapsto-v$ and $\phi:v\mapsto v^3$. Its group is $C_2\times C_m$.
Inversion from interchanging the two parabolics is already in this group.

Suppose first that $m\ne3$. If $\phi(x+\F_3)=x+\F_3$, then
$x^3-x=c\in\F_3$. Relative trace gives $mc=0$, whence $c=0$
and $x\in\F_3$. Thus no nonzero vector is fixed by $\phi$, or by
any of its nonidentity powers. A map $-\phi^k$ fixes no nonzero vector
either: its $m$th power is negation. Consequently the action on
$V\setminus0$ is free, proving the formula.

For $m=3$, $V$ has dimension two over $\F_3$. The induced $\phi$
has order dividing three and is not the identity. To see the latter,
if every $x\in\F_{27}$ satisfied $x^3-x\in\F_3$, the additive
map $x\mapsto x^3-x$ would have image of size at most three, whereas
its kernel has size three and its image has size nine. Thus $\phi$
has one nontrivial unipotent Jordan block on $V$, and its fixed subspace
has dimension one. The two nonidentity powers of $\phi$ each fix two
nonzero vectors. The three maps $-\phi^k$ fix none. Burnside's formula
gives $(8+2+2)/6=2$.
\end{proof}

\section{Source joins and attribution}
Theorem~\ref{four} recovers the $A_5,E_{16}:3$ counts printed in
Table 1, rows 2--4, and Lemmas 5.3--5.4 of the published Part II paper.
Lemma~\ref{subgroups} separately counts subgroup classes, not geometries;
these two counts must not be confused. The same mechanism supplies the
parameter space for the omitted higher-extension subfield covers already
proved in the accepted correction. It does not repeat that acquisition.

Theorem~\ref{three} corrects the inner conjugacy column in Table 5:
the printed $3^{m-1}-1$ becomes $(3^{m-1}-1)/2$, and the printed eight
at $q=27$ becomes four. Both printed automorphism counts remain correct.
The separately fixed 21-page appendix repeats the ordered count in its
proof of Proposition 26, physical page 15. Its calculation counts the
possible second parabolics after choosing a first one; it omits the
additional identification required by the source's unordered definition.
The explicit conjugator $u_{-t}$ carries the unordered pair
$\{H_0,H_t\}$ to $\{H_{-t},H_0\}$. This gives the missing
identification inside $G$, without a computational assumption.

The original numerical formulas recovered here belong to Buekenhout,
De Saedeleer and Leemans. Gaussian counts, finite-field trace, Burnside's
formula and normalizer actions are standard. The contribution of this
note is the common projective parameter description, its extension to
the already established missing covers, and the explicit audit of the
unordered convention. No historical priority beyond this fixed-source
comparison, complete treatment of either original, Lean verification,
new open-conjecture answer or compression measurement is claimed.

\begin{thebibliography}{9}
\bibitem{source} F. Buekenhout, J. De Saedeleer and D. Leemans,
\emph{On the rank two geometries of the groups $\PSL(2,q)$: part II},
Ars Mathematica Contemporanea \textbf{6} (2013), 365--388.
\bibitem{appendix} The same authors, \emph{Part II (APPENDIX)},
the fixed separately catalogued 21-page edition, proof of Proposition 26.
\bibitem{prior} Human Compression Project, \emph{Subfield planes in
projective rank two geometries}; and \emph{Local actions in rank two
over a projective line}, previously accepted components.
\end{thebibliography}
\end{document}
